\documentclass[12pt]{article}
\usepackage[utf8]{inputenc}
\usepackage[T1]{fontenc}
\usepackage{hyperref}
\title{Dharuk Language Resources for K-12 Curriculum Development}
\author{Compiled by Claude Sonnet 5.5 (claude-sonnet-5-5) for Andr\'as Kornai}
\date{}
\begin{document}
\maketitle

\section{Introduction}
Dharuk (ISO 639-3 xdk; Glottocode sydn1236) is usually spelled Dharug. It is the Aboriginal language of the Sydney region and is dormant. Jakelin Troy also calls it the ``Sydney Language''. Wikipedia reports that some effort has gone into reviving a reconstructed form of the language, and that UNESCO classifies it as Critically Endangered. Revival work should follow the lead of the Dharug community.

\section{Primary Educational Resources}
A news item quoted on Wordnik says that staff at Chifley College's Dunheved campus in Sydney have helped revitalise Dharug. I did not verify the current state of any school programme. I saw no URL for the school programme itself.

\section{Linguistic and Reference Materials}
\begin{itemize}
\item Wikipedia, ``Dharug language''. \url{https://en.wikipedia.org/wiki/Dharug_language}
\item Jakelin Troy, \emph{The Sydney Language}, sold through AIATSIS. This is the 1993 history, grammatical comments and extended wordlist. \url{https://shop.aiatsis.gov.au/products/the-sydney-language-dictionary}
\item Jeremy Steele, ``William Dawes's word lists'', held in the SOAS Archives. It is based on Dawes's manuscript, a primary historical source. \url{https://archiveshub.jisc.ac.uk/data/gb102-ms380782}
\end{itemize}

\section{Religious and Community Resources}
No scripture or community organisation was identified with a verified URL.

\section{Digital and Audio Resources}
The Endangered Languages Project lists a sample page for Dharug, which hosts a copy of the Troy wordlist as a Google Doc. \url{https://www.endangeredlanguages.com/lang/6686/samples/11052} No audio was found.

\section{Curriculum Development Context}
The historical record is thin: colonial wordlists and Troy's reconstruction. A curriculum would mostly be a language-awareness and revival programme. It should be built with the Dharug community, and the sources above are starting points only. I did not check whether the Troy wordlist is freely accessible beyond that hosted copy.
\end{document}